
\section{Discussion and Conclusion}
\label{sec:discussion}

We have presented a comprehensive framework for visualizing and exploring 3D integral curves through the Curve Segment Neighborhood Graph (CSNG). Our approach combines efficient community detection algorithms with multi-level visualization techniques, including force-directed layouts and the Adjacency Matrix of Curve Segments (AMCS). Our framework offers a fast, scalable, and qualitatively superior solution for integral curve analysis compared to traditional methods like PCA $\mathbb{K}$-means. The key to its effectiveness is the flexibility to analyze data at both the streamline and segment levels, enabling a hierarchical exploration workflow. Our findings suggest a general strategy: start with the faster, more memory-efficient streamline-level analysis to get an overview, then use the more detailed segment-level analysis on specific regions of interest. This balances performance with analytical depth. Nevertheless, the framework is flexible enough to start with segment-level analysis when complex streamline configurations demand it.

A particularly significant discovery is the effectiveness of low-K KNN configurations ($K \leq 10$) for early-stage feature identification, as they naturally yield sparser graphs that better isolate coherent structures. In the Solar Plume dataset, using $K=10$ led to community \#7 appearing fully separated in the force-directed layout (~\autoref{fig:demo}(f)), capturing the plume’s highly linear core distinct from the surrounding turbulence. This separation highlights how low-K graphs can emphasize physically meaningful regions without prior knowledge of the flow.

% Additional low-K results for feature isolation across more datasets are provided in a detailed supplemental document available on our GitHub repository.

\subsection{Limitations}

While our web-based framework provides exceptional accessibility—-running on any device with a browser and offering a responsive React interface—-this design choice involves significant performance trade-offs. The system suffers from limited multithreading capabilities, where web workers add overhead to runtime, hard memory limits that prevent utilizing full system memory, and the need to avoid UI freezing by delegating expensive tasks to web workers, which introduces data transfer overhead and garbage collection variability. Large operations like community detection or neighbor search can vary by 10-20\% in runtime with no external factors, and these limitations become more pronounced as the size of the data increases.

In addition, a significant challenge in developing our system was fine-tuning the LLM system prompt to understand and adjust the complex interplay of community detection parameters (KNN $K$ values, Louvain resolution) and rendering parameters (opacity thresholds, filtering conditions). We spent considerable time iteratively refining the prompt to enable intelligent parameter recommendations based on dataset characteristics and user goals, significantly reducing the expertise barrier while maintaining flexibility for expert users.

%Really need to clarify R1's misunderstanding that this is not time-depended!
\changed{Our system focuses on analyzing individual time steps or steady-state flow fields. We do not currently support temporal tracking of communities across time steps or ensemble analysis. While our real-time performance could enable frame-by-frame analysis, extending CSNG to encode temporal adjacency relationships would enable automatic analysis across temporal dimensions. Additionally, the interactive split/merge operations introduce subjectivity, as different users may refine communities differently based on their domain knowledge. This is intentional for exploratory analysis and hypothesis generation, but users requiring reproducible results should rely on the automatic community detection baseline. Finally, unlike topology-based feature extraction methods, our neighborhood-based approach does not guarantee identification of all structures of interest—-detected communities depend on neighbor search parameters and resolution settings. We recommend exploring multiple parameter configurations and validating results with domain knowledge or quantitative metrics when available.}

\subsection{Future Work}

While our current implementation shows considerable promise, we have identified several directions for future research. The impact of different neighbor search strategies (RBN versus KNN) on CSNG construction and subsequent analysis requires further investigation. We plan to explore more sophisticated curve-centered neighbor search methods to address the limitations of current approaches and reduce parameter sensitivity. Additionally, we aim to enhance the handling of local patterns and global structures by investigating alternative community detection algorithms and exploring methods for encoding multi-scale flow information into CSNGs.

Our future work will also focus on improving the scalability of our web-based implementation, particularly for processing datasets with a very large number of segments. We plan to explore WebGPU acceleration for computationally intensive operations, though Louvain and $\mathbb{K}$-means remain fundamentally CPU-bound methods. Finally, we will conduct more thorough comparisons between our community detection approach and conventional clustering methods to better understand their relative strengths and limitations across different types of flow datasets. \nguyen{Additionally, our current LLM implementation relies solely on textual descriptions and cannot analyze visual outputs. Future work will explore vision-capable LLMs to enable a closed-loop system where the model can perceive visualization results and iteratively refine parameters based on visual feedback.}
